% TOP_PROBLEM: {"id": 30007433, "title": "Uniformly flat real Littlewood polynomials", "category_id": 8, "category": {"id": 8, "name": "analysis", "display_name": "Analysis", "description": "Limits, continuity, calculus, and function theory.", "slug": "analysis", "order_index": 8, "created_at": "2026-07-31T15:26:25.671Z"}, "status": "open"}
% FIRST_POSTED: 2026-10-06
% ATTEMPT_STATUS: solved
% OPENAI_CLAIM: {"schema_version":1,"source_repo":"https://github.com/openai/math","source_commit":"adc7f1241b42e322a6451854ab7e4b4c146bf78a","release_date":"2026-10-06","match":"full","resolution":"proved","independently_reviewed":false,"second_pass":"agreed","families":[{"family":"076","title":"Real ultraflat Littlewood polynomials and unbounded binary merit factors","subject":"Real and complex analysis","manuscripts":[{"title":"Ultraflat real Littlewood polynomials","pdf_path":"preprints/Ultraflat-real-Littlewood-polynomials-October-5-2026/ultraflat-real-littlewood-polynomials.pdf","readme_path":"preprints/Ultraflat-real-Littlewood-polynomials-October-5-2026/README.md","date":"2026-10-05","theorem_ref":"preprints/Ultraflat-real-Littlewood-polynomials-October-5-2026/build/sections/introduction.tex :: thm:main"}],"lean":{"doc_path":"lean/docs/076.md","comparators":[],"declarations":[],"covers_main_theorem":false},"reasoning_trace":null}]}
\documentclass[11pt]{article}
\usepackage[margin=25mm]{geometry}
\usepackage{fontspec}
\setmainfont{Latin Modern Roman}
\usepackage{amsmath,amssymb,mathtools}
\usepackage{unicode-math}
\setmathfont{Latin Modern Math}
\usepackage{enumitem}
\usepackage{xurl,hyperref,bookmark}
\hypersetup{colorlinks=true,urlcolor=blue}
\setcounter{secnumdepth}{0}
\setlength{\parindent}{0pt}
\setlength{\parskip}{5pt}
\setlength{\emergencystretch}{3em}
\begin{document}
\section*{Uniformly flat real Littlewood polynomials}

\subsection{OpenAI's claim}
\textit{Uniformly flat real sign polynomials}: OpenAI claims that real sign polynomials can have normalized modulus uniformly as close to one as desired once the integer length is large enough. Choices at different lengths are independent of one another.

\subsection{Outline of the argument}
\textit{Uniformly flat real sign polynomials}: A continuous circle function is built with nearly constant modulus, real coefficients of nearly sign size in the retained frequency range, and a small Fourier tail. Auxiliary-torus waves and phase corrections coordinate the upper and lower modulus bounds. Parseval controls the aggregate coefficient defect, allowing discrepancy rounding to preserve uniform flatness.

\subsection{Scope relative to this problem}
\textit{Uniformly flat real sign polynomials}: The primary statement says: For every epsilon>0 and every sufficiently large integer length N, a real sign polynomial has modulus between (1-epsilon)sqrt(N) and (1+epsilon)sqrt(N) everywhere on the circle. The target retains these quantified objects and hypotheses, so its positive assertion is covered.

With the hypotheses and conventions above, the claimed result addresses the whole mathematical assertion on this page.

\subsection{Formal verification}
Family 076 describes Lean coverage for near-minimal upper maxima and convergence of finite-exponent modulus means; the later uniform lower-envelope conclusions are separate.

\textit{Uniformly flat real sign polynomials}: No manuscript-specific Comparator is listed in the pinned inventory.

The documented selection does not cover the entire manuscript assertion relevant here.

This site has not run the Lean development. These links describe OpenAI’s selected formal statements and should be read with their stated scope.

\subsection{Status caveat}
OpenAI reports that an internal model produced the manuscripts. Their arguments have not been independently verified as proofs by this site. Any portion outside the documented formal scope may contain mistakes, and the catalogue label reports the announced mathematical claim.

\subsection{Sources}
\begin{itemize}
\item OpenAI, \emph{Ultraflat real Littlewood polynomials}, 2026-10-05; citation key \texttt{OAI:Ultraflat-real-Littlewood-polynomials-October-5-2026}. \href{https://github.com/openai/math/blob/adc7f1241b42e322a6451854ab7e4b4c146bf78a/preprints/Ultraflat-real-Littlewood-polynomials-October-5-2026/ultraflat-real-littlewood-polynomials.pdf}{PDF (GitHub)}; \href{https://github.com/openai/math/raw/adc7f1241b42e322a6451854ab7e4b4c146bf78a/preprints/Ultraflat-real-Littlewood-polynomials-October-5-2026/ultraflat-real-littlewood-polynomials.pdf}{PDF (direct)}; \href{https://github.com/openai/math/blob/adc7f1241b42e322a6451854ab7e4b4c146bf78a/preprints/Ultraflat-real-Littlewood-polynomials-October-5-2026/README.md}{Manuscript page and citation}; \href{https://github.com/openai/math/blob/main/preprints/Ultraflat-real-Littlewood-polynomials-October-5-2026/ultraflat-real-littlewood-polynomials.pdf}{Latest version}.
\item Family 076: \href{https://github.com/openai/math/blob/adc7f1241b42e322a6451854ab7e4b4c146bf78a/lean/docs/076.md}{Lean scope}.
\item \href{https://github.com/openai/math/blob/adc7f1241b42e322a6451854ab7e4b4c146bf78a/CONTENTS.md}{OpenAI catalogue entry}.
\end{itemize}
\end{document}
